\section*{अध्याय \ref{ch:g4:air-a-mixture-of-gases} --- वायु, गैसों का एक मिश्रण}

\begin{solution}{exo:g4:air-a-mixture-of-gases:1}
एक मिश्रण। उसकी दो मुख्य गैसें नाइट्रोजन और ऑक्सीजन हैं।
\end{solution}

\begin{solution}{exo:g4:air-a-mixture-of-gases:2}
\omterm{def:g4:air-a-mixture-of-gases:air}{वायु} का लगभग चार-पाँचवाँ भाग नाइट्रोजन है, लगभग पाँचवाँ भाग ऑक्सीजन।
\end{solution}

\begin{solution}{exo:g4:air-a-mixture-of-gases:3}
लगभग 78 लीटर नाइट्रोजन और लगभग 21 लीटर ऑक्सीजन।
\end{solution}

\begin{solution}{exo:g4:air-a-mixture-of-gases:4}
बोतल ख़ाली नहीं है: वह \omterm{def:g4:air-a-mixture-of-gases:air}{वायु} से भरी है, और फँसी हुई \omterm{def:g4:air-a-mixture-of-gases:air}{वायु} पानी को अंदर
नहीं आने देती।
\end{solution}

\begin{solution}{exo:g4:air-a-mixture-of-gases:5}
ऑक्सीजन कम है (21 के बजाय 100 में 16 भाग) और कार्बन डाइऑक्साइड ज़्यादा
है (लगभग कुछ नहीं के बजाय 100 में 4 भाग)।
\end{solution}

\begin{solution}{exo:g4:air-a-mixture-of-gases:6}
दोनों में ऑक्सीजन: लौ और हमारा शरीर, दोनों \omterm{def:g4:air-a-mixture-of-gases:air}{वायु} की ऑक्सीजन का
उपयोग करते हैं।
\end{solution}

\begin{solution}{exo:g4:air-a-mixture-of-gases:7}
78 वर्ग नाइट्रोजन हैं; $100 - 78 = 22$ वर्ग नाइट्रोजन नहीं हैं।
\end{solution}

\begin{solution}{exo:g4:air-a-mixture-of-gases:8}
लौ मर्तबान की \omterm{def:g4:air-a-mixture-of-gases:air}{वायु} की ऑक्सीजन की खपत करती है। जब लौ के आस-पास बहुत कम
ऑक्सीजन बचती है, तो वह बुझ जाती है।
\end{solution}

\begin{solution}{exo:g4:air-a-mixture-of-gases:9}
500 लीटर का पाँचवाँ भाग: $500 \div 5 = 100$ लीटर ऑक्सीजन।
\end{solution}

\begin{solution}{exo:g4:air-a-mixture-of-gases:10}
कमरे में मौजूद लोग ऑक्सीजन अंदर लेते हैं और कार्बन डाइऑक्साइड बाहर छोड़ते
हैं। खिड़कियाँ खोलने से ताज़ी \omterm{def:g4:air-a-mixture-of-gases:air}{वायु} वापस आती है, जिसमें ऑक्सीजन ज़्यादा और
कार्बन डाइऑक्साइड कम होती है।
\end{solution}

\begin{solution}{exo:g4:air-a-mixture-of-gases:11}
तिगुनी \omterm{def:g4:air-a-mixture-of-gases:air}{वायु} में तिगुनी ऑक्सीजन होती है: लगभग
$3 \times 10 = 30$ सेकंड।
\end{solution}
